\section*{Chapitre \ref{ch:iv1:linear-algebra} --- Algèbre linéaire}
\markboth{\protect\omTocMark{}{Réponses~: Algèbre linéaire}}{\protect\omTocMark{}{Réponses~: Algèbre linéaire}}

\begin{solution}{iq:iv1:linear-algebra:1}
La matrice $uv^\top$ est de rang 1, et ses valeurs propres sont $v^\top u = 6$, 0 et 0. La matrice de tous les uns a pour valeurs propres $n$ et $0$ ($n - 1$ fois). Avec des zéros sur la diagonale, elles valent $n - 1$ et $-1$ ($n - 1$ fois).

Pour tout vecteur $x$, $uv^\top x = u(v^\top x)$\index{matrices de rang un}, un multiple de $u$. L'image est donc une droite, et $u$ est un vecteur propre~: $uv^\top u = (v^\top u)u$. Tout ce qui est orthogonal à $v$ est envoyé sur 0. La trace vaut $v^\top u$, ce qui sert de contrôle.

La matrice de tous les uns vaut $\mathbf{1}\mathbf{1}^\top$, avec $\mathbf{1}^\top\mathbf{1} = n$. Retrancher l'identité décale chaque valeur propre de $-1$.

\iqwrong{Calculer un polynôme caractéristique. Repérer d'abord la structure~: un décalage par l'identité déplace toutes les valeurs propres, et laisse les vecteurs propres inchangés.}
\end{solution}

\begin{solution}{iq:iv1:linear-algebra:2}
Les valeurs propres sont $2.8$ une fois, avec pour vecteur propre $\mathbf{1}$, et $0.8$ neuf fois, sur tout ce qui est orthogonal à $\mathbf{1}$. Le $\rho$ minimal vaut $-1/9$, et la variance de la moyenne tend vers $\rho$ lorsque le nombre d'actifs devient grand.

On écrit la matrice sous la forme $(1 - \rho)I + \rho\mathbf{1}\mathbf{1}^\top$\index{equicorrelation@équicorrélation}. La matrice de tous les uns a pour valeurs propres $n$ et 0, donc les valeurs propres ici sont $1 - \rho$ et $1 - \rho + n\rho = 1 + (n - 1)\rho$.

Une matrice de corrélation exige que toutes les valeurs propres soient positives ou nulles, donc $\rho \ge -1/(n - 1)$. Pour cette valeur, la somme des actifs est constante.

La moyenne a pour variance $\frac{1}{n^2}\mathbf{1}^\top C\mathbf{1} = \frac{1 + (n - 1)\rho}{n} = \rho + \frac{1 - \rho}{n}$. Ici, elle vaut $0.28$. C'est la limite de la diversification~: la part commune ne disparaît jamais.

\iqwrong{Croire que tout $\rho$ dans $[-1, 1]$ est possible. Avec beaucoup d'actifs, ils ne peuvent pas tous être fortement corrélés négativement entre eux.}
\end{solution}

\begin{solution}{iq:iv1:linear-algebra:3}
Non. Son déterminant vaut $-0.62$, et le portefeuille $(1, -\sqrt{2}, 1)$ a une variance de $4 - 3.6\sqrt{2} \approx -1.09$, ce qui est impossible.

Une matrice de corrélation valide doit être semi-définie positive\index{matrices semi-définies positives}, car $w^\top Cw$ est la variance d'un portefeuille. Tous les blocs 2 par 2 semblent corrects, donc regarder les paires ne permet pas de le détecter. Par symétrie, essayons $w = (1, -a, 1)$. Sa variance vaut $1 + a^2 + 1 - 2a(0.9) - 2a(0.9) = 2 + a^2 - 3.6a$. Elle est minimale pour $a = 1.8$, et déjà négative pour $a = \sqrt{2}$.

Intuitivement, si $X$ est proche de $Y$ et $Y$ proche de $Z$, alors $X$ ne peut pas être décorrélée de $Z$.

\iqwrong{Vérifier que chaque nombre appartient à $[-1, 1]$. Ce sont des conditions nécessaires sur les paires, mais l'exigence porte sur la matrice entière.}
\end{solution}

\begin{solution}{iq:iv1:linear-algebra:4}
Les valeurs propres ne peuvent valoir que 0 ou 1. Le rang est égal à la trace, et la trace de $X(X^\top X)^{-1}X^\top$ vaut $p$.

Si $Pv = \lambda v$, alors $P^2v = \lambda^2v$, donc $\lambda^2 = \lambda$\index{matrices de projection}. Comme $P(P - I) = 0$, le polynôme $x(x - 1)$ a des racines distinctes, donc $P$ est diagonalisable. La trace est le nombre de 1, et le rang est aussi le nombre de 1.

Pour la matrice chapeau, $\operatorname{tr}\bigl(X(X^\top X)^{-1}X^\top\bigr) = \operatorname{tr}\bigl((X^\top X)^{-1}X^\top X\bigr) = \operatorname{tr}(I_p) = p$. La règle de cyclicité fait tout le travail.

\iqwrong{Supposer qu'une projection doit être symétrique. La matrice dont les lignes sont $(1, 1)$ et $(0, 0)$ est idempotente et non symétrique, et son rang et sa trace restent égaux.}
\end{solution}

\begin{solution}{iq:iv1:linear-algebra:5}
Les valeurs propres sont 3 et 1, avec pour vecteurs propres $(1, 1)$ et $(1, -1)$. Ainsi $A^{10}$ a sur la diagonale $(3^{10} + 1)/2 = 29\,525$ et hors de la diagonale $(3^{10} - 1)/2 = 29\,524$, et $A^{-1}$ a sur la diagonale $2/3$ et hors de la diagonale $-1/3$.

La matrice vaut $I + \mathbf{1}\mathbf{1}^\top$, donc ses valeurs propres s'obtiennent immédiatement\index{décomposition spectrale}. Elle est égale à 3 fois une projection plus 1 fois une autre, et le produit de ces deux projections est nul. Ainsi toute fonction de $A$ n'agit que sur les valeurs propres 3 et 1~:
\[
f(A) = \frac{f(3)}{2}\begin{pmatrix}1&1\\1&1\end{pmatrix} + \frac{f(1)}{2}\begin{pmatrix}1&-1\\-1&1\end{pmatrix}.
\]
Avec $f(x) = 1/x$, la diagonale vaut $(\tfrac13 + 1)/2 = \tfrac23$, et les termes hors diagonale valent $(\tfrac13 - 1)/2 = -\tfrac13$. Une multiplication directe donne le même résultat.

\iqwrong{Élever les coefficients à la puissance 10. Les puissances d'une matrice agissent sur ses valeurs propres, et non sur ses coefficients.}
\end{solution}

\begin{solution}{iq:iv1:linear-algebra:6}
La corrélation entre $X$ et $Z$ peut prendre n'importe quelle valeur entre $0.62$ et 1.

Par le déterminant~: une matrice de corrélation 3 par 3 est valide lorsque\index{bornes de corrélation}
\[
1 + 2abc - a^2 - b^2 - c^2 \ge 0,
\]
ce qui, pour $a$ et $b$ fixés, donne $c \in \bigl[ab - \sqrt{(1 - a^2)(1 - b^2)},\ ab + \sqrt{(1 - a^2)(1 - b^2)}\bigr]$. Ici, cela donne $0.81 \pm 0.19$.

Par les angles~: une corrélation est le cosinus de l'angle entre deux variables aléatoires vues comme des vecteurs. $\arccos 0.9 \approx 25.8$ degrés. Les angles entre $X$ et $Z$ peuvent aller de 0 à $2 \cdot 25.8$ degrés, et le cosinus à la borne supérieure vaut $\cos(2\arccos 0.9) = 2 \cdot 0.81 - 1 = 0.62$.

La corrélation n'est pas transitive en général. Avec des corrélations de $0.1$ et $0.1$, la troisième peut prendre n'importe quelle valeur entre $-0.98$ et 1, ce qui ne constitue aucune contrainte réelle. Une contrainte forte n'apparaît que lorsque les corrélations sont élevées.

\iqwrong{Multiplier $0.9 \cdot 0.9 = 0.81$ et donner ce résultat comme réponse. C'est le centre de l'intervalle, et ce serait la réponse si $X$ et $Z$ n'étaient liés que par l'intermédiaire de $Y$. Ce n'est pas nécessairement le cas.}
\end{solution}

\begin{solution}{iq:iv1:linear-algebra:7}
Les valeurs propres de $P$ sont les racines $n$-ièmes de l'unité, les vecteurs propres sont les vecteurs de Fourier $(1, \omega, \omega^2, \dots)$ et le conditionnement vaut 1. La matrice circulante a pour valeurs propres $3$, $1$, $2 + i$ et $2 - i$.

Si $x_k = \omega^k$, avec $\omega^n = 1$, alors $(Px)_k = \omega^{k + 1} = \omega x_k$\index{matrices circulantes}. Comme $P$ est orthogonale, $P^\top P = I$, toutes ses valeurs singulières valent 1, ce qui donne le conditionnement.

Une matrice circulante de première ligne $(c_0, c_1, \dots)$ est le polynôme $c_0I + c_1P + \dots$ en $P$. Elle a donc les mêmes vecteurs propres, et ses valeurs propres sont $c_0 + c_1\omega + \dots$. Ici, $2 + \omega$ avec $\omega \in \{1, i, -1, -i\}$ donne $3$, $2 + i$, $1$ et $2 - i$.

\iqwrong{Affirmer qu'une matrice réelle a des valeurs propres réelles. Seules les matrices symétriques le doivent. Un décalage est une rotation déguisée, et une rotation a des valeurs propres complexes.}
\iqwhy{C'est pourquoi la transformée de Fourier rapide résout les convolutions circulaires~: la base de Fourier diagonalise toutes les matrices circulantes à la fois.}
\end{solution}

\begin{solution}{iq:iv1:linear-algebra:8}
Pour $n = 3$, les valeurs propres sont $2 - \sqrt{2}$, $2$ et $2 + \sqrt{2}$. En général, elles valent $2 - 2\cos\frac{k\pi}{n + 1}$ pour $k = 1, \dots, n$. Le conditionnement croît comme $4(n + 1)^2/\pi^2$, soit environ 4100 pour $n = 100$.

Essayons $x_j = \sin\frac{jk\pi}{n + 1}$, une onde sinusoïdale discrète nulle en $j = 0$ et $j = n + 1$\index{matrice des différences secondes}. Alors
\[
-x_{j - 1} + 2x_j - x_{j + 1} = \Bigl(2 - 2\cos\frac{k\pi}{n + 1}\Bigr)x_j,
\]
par la formule d'addition des sinus. Elle vérifie d'elle-même les lignes de bord.

La plus grande valeur propre est proche de 4, et la plus petite vaut environ $\pi^2/(n + 1)^2$, donc le rapport se comporte comme annoncé. Une grille plus fine donne une matrice moins bien conditionnée. Les résolutions directes restent satisfaisantes pour un système tridiagonal, mais les erreurs d'arrondi croissent avec la grille, et un solveur itératif ralentit sans préconditionneur.

\iqwrong{Croire qu'une grille plus fine donne toujours de meilleures réponses. L'erreur de discrétisation diminue, mais l'erreur d'arrondi augmente, et il existe une taille de grille optimale.}
\end{solution}

\begin{solution}{iq:iv1:linear-algebra:9}
Avec des volatilités unitaires, les composantes sont $(1, 1)/\sqrt{2}$, avec une part de 75~\%, et $(1, -1)/\sqrt{2}$, avec 25~\%. Si la volatilité du premier actif est doublée, la première composante explique environ 86~\%. L'ACP dépend des unités, donc on utilise la matrice de corrélation lorsque les unités ne sont pas comparables.

La première matrice a pour valeurs propres $1 + \rho = 1.5$ et $1 - \rho = 0.5$\index{analyse en composantes principales}, dont la somme vaut 2.

Avec des volatilités de 2 et 1, la matrice de covariance a pour lignes $(4, 1)$ et $(1, 1)$. Ses valeurs propres valent $\frac{5 \pm \sqrt{13}}{2}$, soit $4.30$ et $0.70$. La part expliquée vaut $4.30/5 = 0.86$. La première composante principale pointe maintenant surtout dans la direction du premier actif, à environ 17 degrés de celui-ci, avec $\tan 2\theta = 2/3$. La corrélation n'a pas changé. Seules les unités, elles, ont changé.

Ainsi, l'ACP appliquée à des rendements aux volatilités très différentes fait surtout ressortir les actifs les plus volatils. Il faut décider si l'on souhaite ce résultat. Sinon, il faut normaliser à variance unitaire, ou pondérer par un budget de risque.

\iqwrong{Traiter les composantes comme des propriétés des actifs. Elles dépendent aussi de l'échelle choisie, et un changement de devise ou un facteur d'effet de levier les modifie.}
\end{solution}

\begin{solution}{iq:iv1:linear-algebra:10}
Les valeurs singulières valent $3\sqrt{5} \approx 6.71$ et $\sqrt{5} \approx 2.24$, donc le conditionnement vaut 3. Les valeurs propres valent 3 et 5, ce qui suggérerait un rapport différent.

Les valeurs singulières sont les racines carrées des valeurs propres de $A^\top A$\index{décomposition en valeurs singulières}. Ici $A^\top A$ a pour lignes $(25, 20)$ et $(20, 25)$, dont les valeurs propres valent 45 et 5.

La matrice est triangulaire, donc ses valeurs propres sont sa diagonale~: 3 et 5. Pour une matrice non symétrique, les valeurs propres ne mesurent ni la taille ni le conditionnement, seules les valeurs singulières le font.

Pour une matrice de données centrée $X = U\Sigma V^\top$ à $n$ lignes, la covariance vaut $X^\top X/(n - 1) = V\Sigma^2V^\top/(n - 1)$. Les colonnes de $V$ sont donc les composantes principales, et leurs variances valent $\sigma_i^2/(n - 1)$. Un test sur des données aléatoires le confirme. Il vaut mieux utiliser la SVD sur $X$ qu'un solveur de valeurs propres sur la covariance, car former la covariance élève le conditionnement au carré.

\iqwrong{Prendre le rapport $5/3 \approx 1.67$ comme conditionnement. Cela ne vaut que pour une matrice symétrique.}
\end{solution}

\begin{solution}{iq:iv1:linear-algebra:11}
Les solutions sont $(2, 0)$ et $(1, 1)$. Ainsi, modifier $b$ de $0.0001$ déplace $x$ de 1. Le conditionnement vaut environ $4 \times 10^4$, et l'on perd environ $4.6$ chiffres. Former $A^\top A$ élève le conditionnement au carré, ce qui fait perdre environ $9$ chiffres.

La matrice est presque singulière\index{conditionnement}, puisque ses lignes sont presque égales. Ses valeurs propres valent environ 2 et $5 \times 10^{-5}$, et leur rapport est le conditionnement.

La règle empirique veut que l'erreur relative sur la réponse puisse valoir $\kappa$ fois l'erreur relative sur les données. Si l'on travaille avec 16 chiffres, on peut en perdre environ $\log_{10}\kappa$. Avec les équations normales, $X^\top X$ a un conditionnement de $\kappa^2$, et c'est pourquoi un bon code de régression utilise la QR ou la SVD sur $X$.

\iqwrong{Vérifier le résidu et se satisfaire de sa petitesse. Les deux solutions ont des résidus minuscules par rapport à l'autre $b$. Un petit résidu ne signifie pas une petite erreur lorsque $A$ est mal conditionnée.}
\iqwhy{Deux actifs presque identiques dans un modèle de risque créent exactement ce problème. Une petite variation des moyennes ou des covariances estimées donne alors des pondérations de portefeuille très différentes.}
\end{solution}

\begin{solution}{iq:iv1:linear-algebra:12}
Le rang vaut au plus 249, après soustraction des moyennes. Il existe donc au moins 251 portefeuilles indépendants de variance empirique exactement nulle, vers lesquels un optimiseur de variance minimale envoie ses pondérations.

$S = \frac{1}{T - 1}\sum_t (r_t - \bar{r})(r_t - \bar{r})^\top$ est une somme de $T$ matrices de rang un\index{matrice de covariance empirique}, et les vecteurs centrés s'annulent à la somme, donc son rang vaut au plus $T - 1$. Une simulation avec 250 et 500 le confirme avec 249.

Un portefeuille de variance empirique nulle n'est pas sans risque. Ses pondérations sont ajustées aux particularités de l'échantillon, et son risque réel est ordinaire. L'optimiseur n'admet pas de solution unique et, même avec des contraintes, il se dirige tout droit vers ces directions.

Deux remèdes~: imposer une structure avec un modèle factoriel, $\Sigma = D + BB^\top$, qui est de plein rang avec peu de paramètres. Ou rétrécir~: $\alpha F + (1 - \alpha)S$, avec une cible simple $F$. Les limites de pondération aident aussi, car elles agissent comme un rétrécissement.

\iqwrong{Inverser avec une pseudo-inverse et déclarer le problème résolu. Cela donne un nombre, mais le problème est que les données sont trop peu nombreuses, et aucune algèbre matricielle n'y ajoute quoi que ce soit.}
\end{solution}

\begin{solution}{iq:iv1:linear-algebra:13}
Les nouvelles valeurs propres valent $4.11$, $1.31$ et $0.68$. Le conditionnement passe de 50 à environ $6.0$. Les vecteurs propres et la trace restent les mêmes.

Avec $\mu = 6.1/3 \approx 2.03$\index{estimateurs à rétrécissement}, chaque valeur propre devient $0.7\lambda + 0.3\mu$, car $I$ a tous les vecteurs pour vecteurs propres. Les valeurs propres empiriques sont trop dispersées~: la plus grande est trop élevée et la plus petite trop basse. Cette transformation les rapproche donc tout en conservant la variance totale.

Les pondérations d'un optimiseur utilisent $\Sigma^{-1}$, elles sont donc dominées par les plus petites valeurs propres, qui sont les moins fiables. En les relevant de $0.1$ à $0.68$, on réduit d'environ sept fois le poids accordé à ces directions.

Le $0.3$ peut être choisi par la formule de Ledoit--Wolf, qui estime l'intensité minimisant l'espérance du carré de la distance à la vraie matrice, à partir du bruit de $S$. Ou alors on utilise la validation croisée sur la variance hors échantillon.

\iqwrong{Croire que le rétrécissement fait perdre de l'information parce qu'il modifie les chiffres. Les valeurs empiriques comportent des erreurs qui les écartent toujours davantage, et les rapprocher corrige ce biais.}
\end{solution}

\begin{solution}{iq:iv1:linear-algebra:14}
Les pondérations sont proportionnelles à $\sigma_2^2 - \rho\sigma_1\sigma_2$ et $\sigma_1^2 - \rho\sigma_1\sigma_2$. La seconde s'annule en $\rho = \sigma_1/\sigma_2 = 0.5$. Pour $\rho = 0.8$, les pondérations valent $(4/3, -1/3)$, et la volatilité vaut environ $8.94~\%$.

$w \propto \Sigma^{-1}\mathbf{1}$\index{portefeuille de variance minimale}, ce qui, pour deux actifs, donne la formule ci-dessus. Une recherche sur grille du poids pour $\rho = 0.8$ trouve le même résultat.

Pour $\rho = 0.8$~: la variance vaut $\frac{16}{9}(0.01) + \frac19(0.04) - 2 \cdot \frac43 \cdot \frac13 \cdot 0.016 = 0.008$.

Pourquoi la variance est plus faible~: la position courte dans le second actif couvre le premier. L'effet marginal de l'ajout du second actif à poids nul est sa covariance avec le premier, $0.016$, qui dépasse la variance du premier, $0.01$. Être court sur le second actif réduit donc le risque. Autrement dit, le bêta de l'actif 1 sur l'actif 2 vaut $0.4$, et le meilleur portefeuille détient une couverture partielle contre celui-ci.

Au-dessus de $\rho = 0.5$, l'actif le plus risqué sert de couverture et il est vendu à découvert. En dessous, on le détient en position longue.

\iqwrong{Affirmer qu'un actif plus risqué reçoit toujours une pondération positive plus faible. Pour une corrélation assez élevée, il reçoit une pondération négative, car il est alors utile comme couverture.}
\end{solution}

\begin{solution}{iq:iv1:linear-algebra:15}
La matrice a pour lignes $(0.9, 0.1)$ et $(0.2, 0.8)$, avec des valeurs propres 1 et $0.7$. La loi stationnaire vaut $(2/3, 1/3)$. L'effet de l'état du jour s'amortit d'un facteur $0.7$ chaque jour, soit une demi-vie d'environ $1.9$ jours.

Les lignes somment à 1, donc $\mathbf{1}$ est un vecteur propre à droite de valeur propre 1\index{chaîne de Markov}. L'autre valeur propre découle de la trace~: $1.7 - 1 = 0.7$. La loi stationnaire $\pi$ vérifie $\pi P = \pi$, ce qui donne $0.1\pi_1 = 0.2\pi_2$, donc $\pi_1 = 2\pi_2$.

La différence entre la probabilité d'un lendemain calme partant d'un état calme et celle partant d'un état de stress vaut $0.9 - 0.2 = 0.7$, qui est la valeur propre. Après $k$ jours, elle vaut $0.7^k$, et atteindre la moitié prend $\ln 0.5/\ln 0.7 \approx 1.94$ jours. Une puissance matricielle à 200 pas donne la ligne stationnaire quel que soit le point de départ.

\iqwrong{Croire que la durée d'un état renseigne sur le mélange~: la durée moyenne d'une période de stress est de 5 jours et celle d'une période calme de 10 jours, et ces durées dépassent la demi-vie. Le mélange combine les deux sorties, et la valeur propre les intègre déjà.}
\end{solution}

\begin{solution}{iq:iv1:linear-algebra:16}
La matrice a pour lignes $(1, 1)$ et $(1, 0)$, avec des valeurs propres $\varphi = (1 + \sqrt{5})/2$ et $\psi = (1 - \sqrt{5})/2$. Donc $F_n = (\varphi^n - \psi^n)/\sqrt{5}$, et il faut environ $2\log_2 n$ produits matriciels. Avec des flottants en double précision, la formule arrondie se trompe pour la première fois en $n = 71$.

$(F_{n + 1}, F_n)^\top = M(F_n, F_{n - 1})^\top$, donc $M^n$ a $F_n$ dans sa case hors diagonale\index{puissances matricielles}. Le développement du vecteur initial sur les vecteurs propres donne la formule. Comme $|\psi| < 1$, $F_n$ est l'entier le plus proche de $\varphi^n/\sqrt{5}$.

L'élévation au carré répétée calcule $M^n$ en environ $\log_2 n$ élévations au carré, auxquelles s'ajoutent au plus autant de produits supplémentaires.

Les flottants en double précision disposent d'environ 53 bits. $\varphi^n/\sqrt{5}$ dépasse $2^{53}$ vers $n = 78$, mais les erreurs d'arrondi sur $\varphi$ s'accumulent avant cela. Un test trouve la première valeur fausse en $n = 71$. Il faut donc utiliser les puissances matricielles entières pour obtenir une valeur exacte. Un test calcule ainsi $F_{30} = 832\,040$.

\iqwrong{Affirmer que la forme fermée est toujours meilleure parce qu'elle ne comporte pas de boucle. Elle n'est exacte que dans la mesure de ses entrées en virgule flottante, alors que la méthode entière n'a pas une telle limite.}
\end{solution}

\begin{solution}{iq:iv1:linear-algebra:17}
Le rang vaut 2, et les valeurs propres non nulles valent $5050 \pm \sqrt{100 \cdot 338\,350} \approx 10\,866.8$ et $-766.8$.

Écrivons $A = u\mathbf{1}^\top + \mathbf{1}u^\top$\index{structure de rang faible}, avec $u = (1, \dots, 100)$. Chaque terme est de rang un, et les deux ensemble sont de rang deux. Tout vecteur propre de valeur propre non nulle appartient à l'espace engendré par $u$ et $\mathbf{1}$. Posons $x = au + b\mathbf{1}$. Alors $Ax = u(\mathbf{1}^\top x) + \mathbf{1}(u^\top x)$, ce qui donne la matrice 2 par 2
\[
\begin{pmatrix} \mathbf{1}^\top u & \mathbf{1}^\top\mathbf{1} \\ u^\top u & u^\top\mathbf{1} \end{pmatrix} = \begin{pmatrix} 5050 & 100 \\ 338\,350 & 5050 \end{pmatrix}.
\]
Ses valeurs propres valent $5050 \pm \sqrt{100 \cdot 338\,350}$. Pour vérifier, la somme vaut $10\,100$, ce qui est la trace de $A$~: $\sum 2i = 10\,100$. Un solveur numérique donne le même résultat.

\iqwrong{S'attendre à ce que toutes les valeurs propres soient positives parce que toutes les entrées sont positives. La matrice n'est pas semi-définie positive. Une entrée positive n'implique pas cela.}
\end{solution}

\begin{solution}{iq:iv1:linear-algebra:18}
Les valeurs propres empiriques se situent entre environ $0.086$ et $2.914$. Dans l'échantillon, le portefeuille affiche une variance de $(1 - q)/N$, avec $q = N/T = 0.5$. Hors échantillon, sa variance réelle vaut $1/((1 - q)N)$. La variance optimale vaut donc $1/N$, l'échantillon en rapporte la moitié, et le portefeuille en livre le double, soit un rapport de 4 entre le risque livré et le risque observé.

La loi affirme que les valeurs propres d'un bruit pur s'étalent jusqu'à $(1 \pm \sqrt{q})^2$\index{loi de Marchenko--Pastur}, soit $(1 \pm 0.707)^2$. Une simulation avec 500 actifs et 1000 jours s'approche des deux bornes.

L'optimiseur prend $w \propto S^{-1}\mathbf{1}$. Il retient donc les faibles valeurs propres empiriques, qui sont artificiellement basses, et se charge sur ces directions. Vingt simulations, une par graine, donnent une variance dans l'échantillon de $0.5/N$ et une variance hors échantillon de $2/N$, à comparer à l'optimum réel $1/N$, qui est ici le portefeuille à pondérations égales.

Pour rendre cela utile, regardons les limites. Quand $q \to 0$, on dispose de beaucoup de données, et les trois valeurs concordent. Quand $q \to 1$, le risque hors échantillon explose. En dessous de la borne supérieure de Marchenko--Pastur, les valeurs propres de données réelles ne peuvent pas être distinguées de celles d'un bruit, il faut donc les nettoyer ou utiliser un modèle factoriel.

\iqwrong{Croire que le risque dans l'échantillon est une estimation fidèle du risque auquel on fera face. Un portefeuille optimisé est le plus flatteur pour cet échantillon, il est donc trop bas par construction.}
\end{solution}

\begin{solution}{iq:iv1:linear-algebra:19}
$\Sigma^{-1} = D^{-1} - \dfrac{D^{-1}\beta\beta^\top D^{-1}}{1 + \beta^\top D^{-1}\beta}$. Cela demande $O(N)$ opérations pour un produit matrice-vecteur, au lieu de $O(N^3)$ pour une inversion directe. Les pondérations de variance minimale sont proportionnelles à $d_i^{-1}(1 - c\beta_i)$, donc les titres à bêta élevé reçoivent des pondérations faibles ou négatives.

Preuve\index{formule de Sherman--Morrison}~: il suffit de développer. Posons $x = D^{-1}\beta$ et $s = \beta^\top x$. Alors
\[
(D + \beta\beta^\top)\Bigl(D^{-1} - \frac{xx^\top}{1 + s}\Bigr) = I + \beta x^\top - \frac{\beta x^\top + \beta s x^\top}{1 + s} = I + \beta x^\top - \beta x^\top = I.
\]
Des tests numériques avec 1000 actifs le confirment, et les poids obtenus par la formule rapide coïncident avec ceux d'une résolution directe.

En appliquant la formule à $\mathbf{1}$, on obtient $w \propto D^{-1}\mathbf{1} - x\frac{\mathbf{1}^\top x}{1 + s}$, dont la $i$-ième composante vaut $\frac{1}{d_i}\bigl(1 - \beta_i\frac{\mathbf{1}^\top x}{1 + s}\bigr)$, ce qui est la forme donnée plus haut. Les poids décroissent donc linéairement avec le bêta, lorsque les risques résiduels sont égaux. C'est le biais connu des portefeuilles de variance minimale vers les faibles bêtas~: le portefeuille réduit son exposition au marché en achetant les titres à faible bêta et en vendant à découvert ceux à bêta élevé.

Le coût compte~: avec $k$ facteurs, la version de Woodbury ne demande que l'inversion d'une seule matrice $k \times k$, pour un coût de $O(Nk^2)$. Les modèles factoriels sont donc autant un outil de calcul qu'un outil statistique.

\iqwrong{Former et inverser toute la matrice à chaque fois. À cette taille, cela coûte un milliard d'opérations, et l'on perd la structure qui rend le résultat stable.}
\end{solution}

\begin{solution}{iq:iv1:linear-algebra:20}
La corrélation partielle vaut $1/3$. En général, elle vaut $\rho/(1 + \rho)$.

À partir des résidus~: en régressant sur $Z$, on obtient $X - 0.5Z$ et $Y - 0.5Z$\index{corrélation partielle}. Chacun a pour variance $1 - 0.25 = 0.75$. Leur covariance vaut $0.5 - 0.25 - 0.25 + 0.25 = 0.25$. La corrélation vaut donc $0.25/0.75 = 1/3$. Une simulation le confirme. En général, $\frac{\rho_{XY} - \rho_{XZ}\rho_{YZ}}{\sqrt{(1 - \rho_{XZ}^2)(1 - \rho_{YZ}^2)}}$.

À partir de l'inverse~: posons $P = C^{-1}$. La corrélation partielle vaut alors $-P_{12}/\sqrt{P_{11}P_{22}}$\index{matrice de précision}. Pour cette matrice à corrélations égales, l'inverse a pour diagonale $\frac{1 + \rho}{(1 - \rho)(1 + 2\rho)}$ et pour termes hors diagonale $\frac{-\rho}{(1 - \rho)(1 + 2\rho)}$, ce qui donne $\rho/(1 + \rho)$.

Les termes hors diagonale de l'inverse mesurent donc les liens directs, une fois prises en compte toutes les autres variables. Un terme nul signifie que deux variables sont conditionnellement indépendantes, dans le cas gaussien.

Cas limites~: $\rho = 0$ donne 0, et $\rho \to 1$ donne $\tfrac12$. La valeur n'est pas 1, car lorsque toutes les variables sont presque identiques, les résidus après retrait de $Z$ restent corrélés par la structure commune.

\iqwrong{Prendre la corrélation $0.5$ pour un lien direct. Une partie passe par $Z$. La corrélation partielle retire cette part, et elle peut beaucoup baisser, voire changer de signe.}
\end{solution}

\begin{solution}{iq:iv1:linear-algebra:21}
Multiplier un vecteur de départ par la matrice, normaliser, et répéter. Il faut environ 20 itérations pour une erreur de $10^{-6}$ sur le vecteur, et 10 pour la valeur propre. Si la deuxième valeur propre vaut $9.9$, il faut environ 1375 itérations. La méthode échoue si le vecteur de départ n'a aucune composante selon le vecteur propre dominant.

Écrivons le point de départ $\sum c_iv_i$\index{itération de la puissance}. Après $k$ itérations, il vaut $\lambda_1^k\bigl(c_1v_1 + c_2(\lambda_2/\lambda_1)^kv_2 + \dots\bigr)$. L'erreur décroît donc comme $(\lambda_2/\lambda_1)^k$, et $0.5^{20} \approx 10^{-6}$.

Le quotient de Rayleigh $x^\top Ax$ a une erreur qui dépend du carré de l'erreur du vecteur, pour une matrice symétrique, il demande donc deux fois moins d'itérations. Un calcul sur une matrice 1000 par 1000 confirme ces deux ordres de grandeur.

Pour $9.9/10$, $k = \ln 10^{-6}/\ln 0.99 \approx 1375$. C'est donc l'écart qui compte, et non la taille.

Échec~: orthogonalité exacte à $v_1$, peu probable à partir d'un départ aléatoire, et de toute façon les erreurs d'arrondi corrigent en général cela d'elles-mêmes. Un test qui retire sans cesse la direction dominante montre que l'on reste bloqué sur l'autre direction. Si $\lambda_1 = \lambda_2$, le vecteur propre dominant n'est pas unique~: l'itération converge vers un vecteur quelconque de ce plan, et le choix dépend du point de départ. L'algorithme de Lanczos est beaucoup plus rapide pour un petit écart, en utilisant le même produit matrice-vecteur, et sans former la matrice. Cela compte lorsque la covariance vaut $XX^\top$ ou lorsqu'il s'agit d'un modèle factoriel.

\iqwrong{Donner une réponse qui dépend de la taille de la matrice. Le coût par itération en dépend, à travers le produit matrice-vecteur. Le nombre d'itérations ne dépend que du rapport des valeurs propres.}
\end{solution}

\begin{solution}{iq:iv1:linear-algebra:22}
Si $ABv = \lambda v$ avec $\lambda \ne 0$, alors $w = Bv$ est non nul, et $BAw = BABv = \lambda Bv = \lambda w$. Les valeurs propres non nulles sont donc communes, avec les mêmes multiplicités. Pour les rendements, on forme la matrice $XX^\top$ de taille 250 par 250, et l'on remonte chaque vecteur propre par $v = X^\top u/\sqrt{\lambda}$.

Avec $A = X$ et $B = X^\top$\index{ACP duale}, $XX^\top$ et $X^\top X$ ont les mêmes valeurs propres non nulles. Si $XX^\top u = \lambda u$ et $u$ est de norme 1, posons $v = X^\top u/\sqrt{\lambda}$. Alors $|v|^2 = u^\top XX^\top u/\lambda = 1$, et $X^\top Xv = X^\top XX^\top u/\sqrt{\lambda} = \lambda X^\top u/\sqrt{\lambda} = \lambda v$. Un calcul avec $T = 250$ et $N = 5000$ le confirme. Le problème aux valeurs propres demande environ $(250/5000)^3 \approx 10^{-4}$ du travail.

Pour le déterminant, l'identité $\det(I + AB) = \det(I + BA)$ découle des valeurs propres non nulles communes~: chaque $1 + \lambda$ apparaît dans les deux, et les autres valent 1. Avec $A = u$ et $B = v^\top$, $BA$ est le nombre $v^\top u$, une matrice de taille 1 par 1. Donc $\det(I + uv^\top) = 1 + v^\top u$. Autrement, $uv^\top$ a une seule valeur propre non nulle, $v^\top u$, et les autres valent 0.

\iqwrong{Former d'abord la matrice $5000 \times 5000$ parce que c'est ce que dit la formule. Son rang est au plus 249, si bien que l'essentiel de ce travail porte sur des zéros.}
\end{solution}

\begin{solution}{iq:iv1:linear-algebra:23}
Les valeurs propres sont $1 + 0.9\sqrt{2} \approx 2.273$, $1$ et $1 - 0.9\sqrt{2} \approx -0.273$. Après troncature et remise à l'échelle, $\rho_{12} = \rho_{23} \approx 0.729$ et $\rho_{13} \approx 0.064$. C'est proche, mais ce n'est pas la matrice de corrélation valide la plus proche~: celle-ci a pour valeurs environ $0.74$ et $0.105$.

Mettons la valeur propre négative à 0\index{matrice de corrélation la plus proche}. Soit $v = (\tfrac12, -\tfrac{1}{\sqrt{2}}, \tfrac12)$ son vecteur propre. Alors
\[
A' = A + 0.273\,vv^\top,
\]
dont la diagonale vaut $(1.068, 1.136, 1.068)$. Ses termes hors diagonale valent $0.900 - 0.273 \cdot 0.354 \approx 0.803$ et $0.273 \cdot \tfrac14 \approx 0.068$. En divisant par les racines carrées des éléments diagonaux, on obtient les nouvelles corrélations.

En quel sens est-elle proche~: avant la remise à l'échelle, $A'$ est la matrice semi-définie positive la plus proche au sens de la norme de Frobenius, à une distance de $0.273$, qui est exactement la valeur propre tronquée. La remise à l'échelle l'éloigne ensuite, car elle rétablit la diagonale unité sans tenir compte de la distance.

La matrice de corrélation la plus proche alterne des projections sur le cône des matrices semi-définies positives et sur les matrices à diagonale unité, avec la correction de Dykstra. Un test avec cette méthode donne $0.74$ et $0.105$, avec une distance à la matrice initiale plus petite. Les deux résultats diffèrent peu dans de nombreux usages, et la troncature est simple et rapide.

\iqwrong{Croire que la remise à l'échelle après troncature donne la matrice la plus proche, ou que l'on peut corriger le problème en ajustant le pire nombre isolé. Toute réparation modifie plusieurs termes à la fois.}
\end{solution}

\begin{solution}{iq:iv1:linear-algebra:24}
$E[\det(A)^2] = n!$. Pour $n = 3$ avec des coefficients $\pm 1$, il y a $2^9 = 512$ matrices, et la moyenne de $\det^2$ vaut 6.

Développons le déterminant selon les permutations\index{déterminants aléatoires}~:
\[
\det(A) = \sum_\sigma \operatorname{sgn}(\sigma)\prod_i a_{i\sigma(i)}.
\]
En élevant au carré, on obtient une double somme sur $\sigma$ et $\tau$. Si $\sigma \ne \tau$, il existe une ligne $i$ où $\sigma(i) \ne \tau(i)$, si bien que $a_{i\sigma(i)}$ apparaît seul dans ce produit (chaque coefficient n'apparaît qu'une fois dans chacun des deux produits), et son espérance est nulle. Si $\sigma = \tau$, le produit vaut $\prod a_{i\sigma(i)}^2$, dont l'espérance vaut 1. Seuls les $n!$ termes diagonaux subsistent donc.

L'énumération des 512 matrices donne exactement 6. Des simulations avec des coefficients gaussiens pour $n = 2$, 4 et 5 concordent. La formule ne demande donc que l'indépendance, une espérance nulle et une variance égale à 1.

La taille quadratique moyenne est donc $\sqrt{n!}$, ce qui fixe le taux de croissance. Une valeur typique est plus petite~: avec des coefficients gaussiens, la moyenne géométrique est inférieure d'un facteur d'environ $1.7\sqrt{n}$, peu de chose à côté de la factorielle. Pour les matrices à coefficients $\pm 1$, la valeur maximale possible est au plus $n^{n/2}$, par l'inégalité de Hadamard\index{inégalité de Hadamard}, car le déterminant est le volume d'un parallélépipède dont les côtés ont pour longueur $\sqrt{n}$. Pour $n = 3$, la borne vaut $5.2$, et l'énumération montre que le maximum réel est 4. Le rapport entre $\sqrt{n!}$ et $n^{n/2}$ décroît comme $e^{-n/2}$, d'après la formule de Stirling. Les vecteurs aléatoires sont donc très loin d'être orthogonaux au sens du volume, même si chaque couple est presque orthogonal.

\iqwrong{Dire que l'espérance du déterminant vaut 0, donc que sa taille typique est petite. Certes, l'espérance vaut 0, puisque chaque terme du développement est un produit de facteurs indépendants d'espérance nulle. Le moment d'ordre 2 en donne la taille, et elle est énorme.}
\end{solution}

\begin{solution}{iq:iv1:linear-algebra:25}
Écrivons $A = \sum_k \lambda_ku_ku_k^\top$ et $B = \sum_l \mu_lv_lv_l^\top$, avec tous les coefficients positifs ou nuls. Alors $A \circ B = \sum_{k, l}\lambda_k\mu_l(u_k \circ v_l)(u_k \circ v_l)^\top$, qui est une somme de matrices semi-définies positives\index{théorème du produit de Schur}.

Vérifions l'étape intermédiaire~: le terme $(i, j)$ de $A \circ B$ vaut $\bigl(\sum_k\lambda_ku_{ki}u_{kj}\bigr)\bigl(\sum_l\mu_lv_{li}v_{lj}\bigr)$, ce qui est égal à $\sum_{k, l}\lambda_k\mu_l(u_{ki}v_{li})(u_{kj}v_{lj})$. Des tests aléatoires sur des matrices de faible rang ne font jamais apparaître de valeur propre négative, au-delà des erreurs d'arrondi.

Pour l'atténuation, $e^{-\lambda|i - j|}$ est la matrice de corrélation d'un processus AR(1), et elle est donc définie positive. Le produit avec une matrice de corrélation est donc semi-définie positive. La diagonale reste égale à 1, car les deux diagonales valent 1. C'est une matrice de corrélation valide\index{atténuation de covariance}. Un test avec $n = 8$ le confirme.

Pour la borne, posons $m = \lambda_{\min}(A)$. Alors $A - mI$ est semi-définie positive, donc $(A - mI) \circ B$ l'est aussi. Cela se développe en $A \circ B - m\operatorname{diag}(B)$. Ainsi
\[
x^\top(A \circ B)x \ge m\sum_i B_{ii}x_i^2 \ge m\min_iB_{ii}|x|^2.
\]
Des tests aléatoires confirment aussi cette inégalité.

C'est pourquoi une atténuation douce est sûre, et une coupure brutale ne l'est pas. Annuler à zéro tous les termes au-delà d'une bande revient à faire un produit terme à terme avec une matrice qui n'est pas semi-définie positive, et un test montre que cela peut produire une valeur propre négative~: la matrice à un facteur avec $\rho = 0.9$ et une bande de largeur 1 en donne une.

\iqwrong{Essayer d'utiliser directement les valeurs propres des deux matrices, comme pour un produit matriciel. Les valeurs propres de $A \circ B$ ne sont pas des produits de valeurs propres. C'est la décomposition en termes de rang un qui permet de conclure.}
\end{solution}
